\chapter{Mission Definitions}
\label{chapter: misison definitions}

To aid with understanding and provide a clear responsibility division, this chapter provides a brief outline of definitions of keywords that will be used throughout the report.

The \textbf{VISTA Mission} encompasses the design, development, manufacturing, qualification, operation and disposal of the \textbf{VISTA Vehicle}. For a successful mission completion, many other components are necessary. They are defined below, but the detailed design of some is outside the scope of the project. Where necessary, the interfaces and contact points are considered.

\begin{itemize}
    \setlength{\itemsep}{1pt}   % space between items
    \setlength{\parskip}{1pt}   % space between paragraphs within items
    \setlength{\parsep}{1pt}    % space between paragraphs
    \item \textbf{Ground Systems:} The entire communications and operational architecture on Earth necessary to orchestrate and manage the functioning of all other mission components.
    \item \textbf{Launcher:} The vehicle that delivers the \textbf{Cruise Stage} from ground into Earth orbit.
    \item \textbf{Cruise Stage:} Manages the orbital cruise from Earth and towards Venus. Monitors and maintains the \textbf{Orbiter} and the \textbf{VISTA Vehicle}. Delivers the \textbf{Orbiter} to a Venus orbit and the \textbf{VISTA Vehicle} to a Venus entry trajectory.
    \item \textbf{Orbiter:} Remains in orbit around Venus for the nominal mission duration and provides a communication link between the \textbf{VISTA Vehicle} and the \textbf{Ground Systems}. Can also conduct other tasks as determined by its design.
    \item \textbf{VISTA Vehicle:} The vehicle whose design is the concern of this work. It consists of the \textbf{Aerobot} and the \textbf{Deployment Stage}.
    \item \textbf{Deployment Stage:} Also referred to as the deployment subsystem, it protects and manages the \textbf{Aerobot}, starting at separation from the \textbf{Cruise Stage}, Venus atmospheric entry, until operational readiness of the \textbf{Aerobot}.
    \item \textbf{Aerobot:} The entire vehicle that floats in Venusian atmosphere for the nominal mission duration. It consists of a \textbf{Gondola} suspended from a \textbf{Balloon}.
    \item \textbf{Balloon:} The envelope containing lifting gas that provides buoyancy for the \textbf{Aerobot}. It interfaces and connects with the \textbf{Gondola}.  
    \item \textbf{Gondola:} The structure suspended beneath the \textbf{Balloon}. It contains the \textbf{Bus} and the \textbf{Instruments}.
    \item \textbf{Bus:} Provides all resources necessary for the \textbf{Instruments} to function correctly for the nominal mission duration. Contains the \textbf{ISRU Subsystem}, among other subsystems.
    \item \textbf{ISRU Subsystem:} In situ resource utilization subsystem. A key part of the \textbf{Aerobot} that is tasked with extracting atmospheric nitrogen to aid with countering lifting gas leakage of the \textbf{Balloon}. The design of the \textbf{ISRU Subsystem} is of particular interest in this work.
    \item \textbf{Instruments:} Also sometimes referred to as "payload" or "payload instruments". The set of instruments composing the scientific suite of the \textbf{Aerobot}, that provides the desired research output over the nominal mission duration.
\end{itemize}


